\section{Backend} 
\label{sec:backend}

El backend está construido con diferentes tecnologías de desarrollo web, en esta sección se describen las herramientas que forman la capa de servidor y persistencia de datos de la aplicación.


\subsection{Python}
\label{subsec:python}

Python es un lenguaje de programación interpretado, de alto nivel y útil para diferentes tipos de tareas. Creado por Guido van Rossum a principios de los años 90~\cite{python_history}, Python ha evolucionado hasta convertirse en uno de los lenguajes más utilizados del mundo, con un 57,9\% de uso entre los desarrolladores en 2025~\cite{stackoverflow_survey_2025}, especialmente en el ámbito
del desarrollo web, el análisis de datos y la inteligencia artificial. Siendo uno de los lenguajes de programación más versátiles.

La sintaxis clara permite escribir programas concisos para combinar lógica del servidor, procesamiento de APIs y generación dinámica de código. Python cuenta con un ecosistema de bibliotecas muy completo, distribuido a través de PyPI (\emph{Python Package Index}), del que WebBuilder hace uso directo: \texttt{requests} para la descarga de datos desde APIs externas, \texttt{xmltodict} y \texttt{defusedxml} para el parseo seguro de XML y \texttt{zipfile} para el empaquetado de los proyectos generados.


En WebBuilder, Python actúa como columna vertebral de la lógica del servidor, gestionando la ingestión y validación de datos, la construcción y envío de prompts al LLM y la generación del código de los proyectos. Los detalles técnicos de cada función se describen en el Capítulo~\ref{chap:diseno}.

La versión utilizada es Python~3, aprovechando características modernas del lenguaje
como las \emph{f-strings} (cadenas de texto con expresiones Python embebidas, usadas en la construcción de prompts para el LLM).

\subsection{Django}
\label{subsec:django}

Django es un framework de desarrollo web de código abierto para Python, que promueve 
el desarrollo rápido y escalable, así como un diseño limpio~\cite{django_docs}. 
Fue publicado en 2005 y sigue el principio \emph{``batteries included''}: incluye 
de serie todo lo necesario para construir aplicaciones web completas, desde el sistema 
de enrutado de URLs hasta el ORM (\emph{Object-Relational Mapper}, capa de software que permite trabajar con la base de datos usando objetos Python en lugar de escribir SQL), el motor de 
plantillas, el sistema de autenticación y el panel de administración.

Permite configurar todo lo necesario, Django incluye un sistema de gestión de base de datos SQLite integrado que permite definir las tablas y sus campos directamente en Python, sin necesidad de escribir SQL. Cada tabla de la base de datos se representa como una clase, y Django se encarga de traducir automáticamente las operaciones sobre esos objetos (consultas, inserciones, actualizaciones) a las instrucciones SQL correspondientes. Esto facilita el desarrollo inicial, y también la evolución del modelo a través de migraciones automáticas, que registran cada cambio de forma controlada, siendo capaz de saber que se cambio o modifico en cada momento, teniendo más control sobre el desarrollo.

La lógica de vistas y de URLs permite organizar la lógica de forma clara y separada, con funciones clave que responden peticiones HTTP, separando el procesamiento de datos de la presentación. El motor de plantillas permite generar HTML dinámico con una sintaxis sencilla, con soporte para herencia de plantillas, lo que favorece la reutilización de componentes visuales.

En WebBuilder, Django gestiona el ciclo de vida completo de cada análisis: desde la recepción y validación de la URL hasta la persistencia de resultados, la invocación del módulo LLM y el servicio de las vistas del asistente. Su funcionamiento detallado se describe en el Capítulo~\ref{chap:diseno}.

La versión empleada es Django~5.1, que introduce mejoras en el ORM y en el sistema de migraciones respecto a versiones anteriores.


\subsection{Base de datos: SQLite y PostgreSQL}
\label{subsec:base-datos}

SQLite es un sistema de base de datos ligero que almacena toda la información en un único fichero local, sin necesidad de instalar ni configurar un servidor externo~\cite{sqlite_docs}. Esta simplicidad lo convierte en la opción por defecto de Django para entornos de desarrollo, aunque presenta limitaciones en escrituras concurrentes al utilizar bloqueos sobre el fichero de datos.

PostgreSQL es un sistema gestor de bases de datos relacional de código abierto, robusto y ampliamente utilizado en entornos de producción~\cite{postgresql_docs}. Ofrece soporte completo para transacciones ACID, concurrencia real mediante MVCC (\emph{Multi-Version Concurrency Control}) y un rico conjunto de tipos de datos, lo que lo convierte en una opción adecuada cuando se requiere mayor robustez y escalabilidad.

En WebBuilder se utilizaron ambas tecnologías de forma progresiva: SQLite durante las fases iniciales de desarrollo por su sencillez, y PostgreSQL en las fases finales para soportar un entorno de producción más exigente. El detalle de la migración de SQLite a PostgreSQL y los motivos técnicos que la justificaron se describen en la Sección~\ref{subsec:migracion-postgresql}.

